\chapter{Eliminatorias y repechaje}
\label{cap:eliminatorias-repechaje}

\begin{procedimiento}{Objetivo del capitulo}{Operador del torneo}
Explicar como registrar ganadores, guardar clasificados multiples, abrir una fase recomendada, evaluar repechaje y crear fases manuales cuando la interfaz lo permita.
\end{procedimiento}

Las eliminatorias y fases posteriores solo deben operarse cuando la fase anterior ya quedo guardada. Si la pantalla solicita completar resultados antes de avanzar, vuelva a la fase anterior y cierre las selecciones pendientes.

\section{Tipos de operacion}

\begin{tabularx}{\linewidth}{>{\bfseries}p{3.5cm}X}
\toprule
Operacion & Uso en la interfaz\\
\midrule
Eliminatoria con ganador & Cada batalla requiere seleccionar un ganador unico con \botoninterfaz{Marcar ganador} y guardar con \botoninterfaz{Guardar este ganador}.\\
Batalla multi-clasificador & La batalla permite seleccionar varios clasificados; la tarjeta muestra \campointerfaz{Clasifican n} y se guarda con \botoninterfaz{Guardar clasificados}.\\
Repechaje & Se ofrece desde el modal de avance cuando existen candidatos suficientes para volver a competir antes de la fase recomendada.\\
Fase manual & Permite preparar una fase personalizada con participantes vivos o eliminados disponibles, distribucion y clasificados por grupo.\\
\bottomrule
\end{tabularx}

\section{Registrar ganador de batalla}

\figuraManual{capturas/finales/MU-026-batalla-ganador.png}{Batallas eliminatorias con ganador seleccionado, participante perdedor y panel lateral de avance de ronda.}{fig:bloque3-mu-026}

\begin{procedimiento}{Guardar ganador de una batalla}{Operador del torneo}
Use este procedimiento en fases donde cada batalla tiene un unico ganador.
\end{procedimiento}

\begin{precondiciones}
\begin{itemize}
    \item La fase esta materializada y aparece en la navegacion.
    \item La batalla tiene participantes visibles.
    \item El resultado oficial ya fue confirmado por mesa o jueces.
\end{itemize}
\end{precondiciones}

\paso{1}{Abra la pestana de la fase correspondiente, por ejemplo una semifinal o batalla de clasificados.}
\paso{2}{Ubique la tarjeta de la batalla.}
\paso{3}{Revise los participantes y el estado visible de la batalla.}
\paso{4}{Presione \botoninterfaz{Marcar ganador} en el equipo que gano.}
\paso{5}{Verifique que el ganador quede resaltado y que los demas queden como salida de ronda.}
\paso{6}{Presione \botoninterfaz{Guardar este ganador}.}
\paso{7}{Confirme el dialogo del navegador si aparece.}
\paso{8}{Revise el mensaje de resultado y el panel lateral \campointerfaz{Ganadores}.}

\begin{resultadoesperado}
La plataforma muestra un mensaje como \campointerfaz{Se registro el ganador de Semifinal demo 1.} y conserva la seleccion al recargar.
\end{resultadoesperado}

\begin{fallas}
Si intenta guardar sin elegir ganador, la interfaz puede mostrar \campointerfaz{selecciona un ganador antes de guardar}. Si el cambio requiere confirmacion por afectar fases posteriores, consulte el procedimiento de correccion de resultados del capitulo correspondiente.
\end{fallas}

\section{Guardar clasificados multiples}

\figuraManual{capturas/finales/MU-027-batalla-clasificados.png}{Batalla multi-clasificador con contador de seleccion, clasificados resaltados y boton de guardado.}{fig:bloque3-mu-027}

\begin{procedimiento}{Guardar clasificados de una batalla multiple}{Operador del torneo}
Use este procedimiento cuando una batalla indica que clasifican varios participantes.
\end{procedimiento}

\paso{1}{Abra la fase donde aparece la batalla multi-clasificador.}
\paso{2}{Revise la etiqueta \campointerfaz{Clasifican n}. La cantidad depende de la fase y no debe asumirse fija.}
\paso{3}{Marque cada equipo clasificado con \botoninterfaz{Clasificar}.}
\paso{4}{Revise el contador \campointerfaz{Seleccionados: x / n}.}
\paso{5}{Corrija la seleccion antes de guardar si el contador no coincide.}
\paso{6}{Presione \botoninterfaz{Guardar clasificados}.}
\paso{7}{Verifique que los clasificados sigan resaltados y aparezcan en el panel lateral de ganadores.}

\begin{resultadoesperado}
La plataforma muestra un mensaje como \campointerfaz{Se guardaron los clasificados de Batalla multi-clasificador demo.}
\end{resultadoesperado}

\begin{fallas}
Si selecciona menos o mas participantes que los requeridos, la interfaz informa la cantidad esperada y la cantidad actual. Ajuste la seleccion y vuelva a guardar.
\end{fallas}

\section{Guardado general en fases posteriores}

El boton \botoninterfaz{Guardar todos los resultados} tambien aparece en fases eliminatorias. Funciona sobre formularios modificados visibles: batallas con ganador cambiado o batallas con clasificados multiples editados.

\begin{procedimiento}{Guardar varias batallas visibles}{Operador del torneo}
Use este recorrido cuando haya cambios revisados en mas de una batalla.
\end{procedimiento}

\paso{1}{Revise cada batalla visible.}
\paso{2}{Confirme que todas las batallas modificadas cumplen su validacion: ganador unico o cantidad exacta de clasificados.}
\paso{3}{Presione \botoninterfaz{Guardar todos los resultados}.}
\paso{4}{Confirme el dialogo \campointerfaz{Seguro que deseas guardar estas selecciones? Podras editarlas despues, pero verifica bien antes de continuar.}}
\paso{5}{Revise el mensaje final con las selecciones guardadas.}
\paso{6}{Si una operacion falla parcialmente, revise el mensaje, recargue la fase y confirme cuales batallas quedaron guardadas.}

\begin{advertencia}
El guardado general no reemplaza la revision de mesa. Use esta accion solo cuando todas las selecciones visibles ya fueron verificadas.
\end{advertencia}

\section{Abrir una fase progresiva}

\figuraManual{capturas/finales/MU-021-modal-siguiente-fase.png}{Modal para pasar a la siguiente fase con recomendacion principal, metricas y alternativas.}{fig:bloque3-mu-021}

\begin{procedimiento}{Revisar el modal de avance}{Operador del torneo}
Use este procedimiento cuando el boton \botoninterfaz{Pasar a la siguiente fase} este habilitado.
\end{procedimiento}

\paso{1}{Revise que todos los resultados de la fase actual esten guardados.}
\paso{2}{Presione \botoninterfaz{Pasar a la siguiente fase}.}
\paso{3}{Lea la recomendacion principal, participantes vivos, eliminados recientes, repechables y formato sugerido.}
\paso{4}{Revise la seccion de \campointerfaz{Alternativas} y \campointerfaz{Advertencias}.}
\paso{5}{Use \botoninterfaz{Cancelar} o cierre el modal si necesita volver a revisar resultados.}

\begin{resultadoesperado}
El operador entiende que el modal es una vista previa: aun no se crea una fase ni se modifican datos hasta confirmar una opcion.
\end{resultadoesperado}

\figuraManual{capturas/finales/MU-022-crear-fase-recomendada.png}{Modal de avance con marcadores editoriales sobre la decision de fase y la zona de revision.}{fig:bloque3-mu-022}

En \cref{fig:bloque3-mu-022}, el marcador 1 senala el area de decision de avance y el marcador 2 encierra la zona que debe revisarse antes de usar la recomendacion.

\begin{procedimiento}{Crear la fase recomendada}{Operador del torneo}
Use esta accion cuando la recomendacion coincide con la decision de la organizacion.
\end{procedimiento}

\paso{1}{Dentro del modal, presione \botoninterfaz{Usar recomendacion del sistema}.}
\paso{2}{Revise la vista previa de creacion: participantes, formato, estructura y unidades.}
\paso{3}{Si aparece decision de repechaje, elija continuar sin repechaje solo si la organizacion asi lo define.}
\paso{4}{Presione \botoninterfaz{Confirmar creacion} o \botoninterfaz{Crear final y tercer lugar}, segun el texto visible.}
\paso{5}{Confirme el dialogo del navegador si aparece.}
\paso{6}{Verifique que la nueva fase aparece en la navegacion y queda disponible para operar.}

\begin{advertencia}
No cree una fase recomendada con resultados incompletos. Si la interfaz muestra \campointerfaz{Completa y guarda todos los resultados antes de pasar a la siguiente fase.}, vuelva a la fase actual y cierre los resultados pendientes.
\end{advertencia}

\section{Repechaje}

\figuraManual{capturas/finales/MU-023-repechaje.png}{Modal de avance con area de evaluacion para repechaje y advertencias de revision.}{fig:bloque3-mu-023}

\begin{procedimiento}{Evaluar o abrir repechaje}{Operador del torneo}
Use este procedimiento solo cuando el modal muestra candidatos suficientes para repechaje.
\end{procedimiento}

\paso{1}{Abra \botoninterfaz{Pasar a la siguiente fase} despues de cerrar la fase actual.}
\paso{2}{Revise el conteo \campointerfaz{Repechables}.}
\paso{3}{Presione \botoninterfaz{Evaluar/Abrir repechaje} si la opcion esta habilitada.}
\paso{4}{Revise candidatos, cupos de retorno, estado y creacion real.}
\paso{5}{Si la organizacion decide usar repechaje, presione \botoninterfaz{Confirmar repechaje}.}
\paso{6}{Confirme el dialogo del navegador: la fase real de repechaje se creara con los candidatos actuales y no se creara todavia la fase principal.}
\paso{7}{Verifique que el repechaje aparece como fase operable y registre sus resultados como una batalla normal o multi-clasificador segun la pantalla.}

\begin{operacionsensible}
El repechaje reincorpora participantes eliminados o recientes candidatos al flujo. Antes de confirmarlo, revise que los candidatos sean correctos y que la fase anterior este completa. Evite modificar resultados anteriores sin revisar consecuencias.
\end{operacionsensible}

\section{Fase manual}

\figuraManual{capturas/finales/MU-024-asistente-manual.png}{Modal de avance con opcion de configurar manualmente una fase personalizada.}{fig:bloque3-mu-024}

\begin{procedimiento}{Generar propuesta de fase manual}{Operador del torneo}
Use una fase manual cuando la recomendacion no cubre la decision operativa o cuando la organizacion necesita una distribucion personalizada.
\end{procedimiento}

\paso{1}{Abra el modal de avance.}
\paso{2}{Presione \botoninterfaz{Configurar manualmente}.}
\paso{3}{Elija \campointerfaz{Fase normal} o \campointerfaz{Repechaje}, segun el origen de participantes.}
\paso{4}{Seleccione una distribucion o use \campointerfaz{Grupos balanceados personalizados}.}
\paso{5}{Defina \campointerfaz{Cantidad de grupos} cuando el modo personalizado lo solicite.}
\paso{6}{Defina \campointerfaz{Clasificados por grupo}.}
\paso{7}{Revise \campointerfaz{Nombre de fase}.}
\paso{8}{Revise el resumen: participantes disponibles, grupos, clasificados y eliminados previstos.}
\paso{9}{Presione \botoninterfaz{Generar propuesta}.}

\figuraManual{capturas/finales/MU-025-propuesta-manual.png}{Recorte de la vista de propuesta manual dentro del modal de avance.}{fig:bloque3-mu-025}

\begin{procedimiento}{Confirmar una propuesta manual}{Operador del torneo}
Use esta accion solo despues de revisar la propuesta generada.
\end{procedimiento}

\paso{1}{Revise participantes asignados, clasificados previstos y eliminados previstos.}
\paso{2}{Si la propuesta muestra grupos, revise cada participante y su grupo destino.}
\paso{3}{Use los selectores \campointerfaz{Mover a} si necesita reasignar participantes dentro de la propuesta.}
\paso{4}{Presione \botoninterfaz{Regenerar grupos} si la propuesta no es aceptable.}
\paso{5}{Presione \botoninterfaz{Confirmar y crear fase} si la propuesta fue aprobada.}
\paso{6}{Use \botoninterfaz{Cancelar} o \botoninterfaz{Cancelar propuesta} si no debe crearse la fase.}
\paso{7}{Verifique que la nueva fase aparece en la navegacion y que sus participantes coinciden con la propuesta.}

\begin{operacionsensible}
Una fase manual puede crear cruces distintos a la recomendacion. Debe usarla un operador autorizado y con validacion de la organizacion. Si la fase no coincide con resultados anteriores, cancele y revise antes de confirmar.
\end{operacionsensible}

\section{Estados visibles en batallas}

\begin{tabularx}{\linewidth}{>{\bfseries}p{3.8cm}X}
\toprule
Estado visual & Interpretacion operativa\\
\midrule
Pendiente & La batalla aun no tiene ganador o clasificados guardados.\\
Ganador & El participante fue seleccionado como ganador unico de la batalla.\\
Clasificado & El participante fue seleccionado en una batalla multi-clasificador.\\
Perdedor o eliminado & La seleccion ya existe y el participante no avanza desde esa batalla.\\
Finalizado & La batalla ya tiene resultado guardado.\\
Bloqueado o incompleto & La fase no puede avanzar porque faltan resultados o candidatos suficientes.\\
\bottomrule
\end{tabularx}

\section{Checklist del capitulo}

\begin{itemize}
    \item Batallas completas y con cantidad correcta de resultados.
    \item Ganadores o clasificados revisados antes de guardar.
    \item Repechaje evaluado solo si la interfaz lo ofrece.
    \item Fase siguiente habilitada despues de guardar todo.
    \item Cambios verificados al recargar la fase.
\end{itemize}
